\begin{lemma}[Marks and localization in fitted frames]
Let $F$ be a finite-dimensional real inner-product space, let
$1\le C_0<\infty$, let $\kappa>0$ and $H\in\R$, and let $v\in\Aclass_{m,C_0}$
in the sense of \noderef{IsFitClass} be centered at $z_*=(0,1)$ with constant
$C_0$ in the sense of \noderef{IsCenteredAt} and satisfy the horizontal Jensen
bound with constant $\kappa$ at every time in the sense of
\noderef{HasJensenBound}. Let $W$ be the contact envelope of
\noderef{ContactEnv} formed from $v$ and $\kappa$, and assume the six standing
properties of the extremal contact pair: $W$ is continuous on $\{t>0\}$ with
subquadratic spatial growth on compact subintervals of $(0,\infty)$ in the
sense of \noderef{SubquadraticGrowth}; $W$ is a viscosity subsolution of the
limiting equation there in the sense of \noderef{IsLimitSubOn};
$v-\kappa^2t/2\le W\le v$ for $t>0$ and $W(z_*)=v(z_*)=0$;
$x\mapsto W(x,t)-|x|^2/(2t)$ is concave for every $t>0$; the barycentric gap of
\noderef{BaryEnv} is $\kappa^2t/2$ at every contact; and the envelope
Hamiltonian of \noderef{EnvHamiltonian} satisfies $H_v(z_*)=H$ with
$\Senv_\ell^+W(z_*)\to H$ as $\ell\downarrow0$, the forward envelope slope being
that of \noderef{EnvSlopeP}. Let $C_*>0$. Write $\delta_j=2^{-j}$ and
$L(\Theta)=1+\log\Theta$.

Then there are an integer $T\ge8$ and reals $R\ge2$, $K\ge2$ and
$C_\Gamma>0$ with the following property. Let $N\ge1$, let
$\epsilon>0$ and let $(\Omega,\mathcal M,\mu)$, $(\mathcal G_n)$, $\xi$, $z$,
$z^+$, $z^-$, $Q_0$, $\tau$ be a stopped contact walk for
$(\kappa,H,T,R,v,N,\epsilon)$ in the sense of \noderef{IsContactWalk}. Write
$z_n=(x_n,t_n)$, $B=R+T+2$,
\[
  Q_+(n)=\frac{x^+_n-x_n}{\epsilon},\qquad
  Q_-(n)=\frac{x_n-x^-_n}{\epsilon},\qquad
  \mathcal E_n=|Q_+(n)-Q_0(n)|^2+|Q_-(n)-Q_0(n)|^2,
\]
and assume the error budget
\begin{equation}\label{eq:ml-contact-error-budget-hyp}
  \mathbb E\sum_{n<\tau}\mathcal E_n\le C_*L(B)^4 .
\end{equation}
Let $(\sigma_j,p_j,a_j)$ be marks with fits for
$(C_0,v,N,\mathcal G,z,\tau)$ in the sense of \noderef{IsMarkData}. Then for
every $j$ with $1\le j\le N$, writing $\Gamma_j=C_\Gamma(L(B)+j)$:
\begin{enumerate}
\item[(i)] $|p_j|+|a_j|\le\Gamma_j$ pointwise, and $\Gamma_j\ge1$;
\item[(ii)] the localization event
\[
  \mathcal L_j=\Bigl\{t_\tau=\delta_N\ \text{and}\
     \bigl|x_\tau-x_{\sigma_j}-(t_\tau-\delta_j)p_j\bigr|\le K\delta_j\Bigr\}
\]
belongs to $\mathcal M$ and, uniformly in $j$ and $N$,
\begin{equation}\label{eq:contact-finite-localization}
  \mathcal L_j\in\mathcal M,\qquad
  \mathbb P\left(t_\tau=\delta_N\quad\text{and}\quad
     \bigl|x_\tau-x_{\sigma_j}-(t_\tau-\delta_j)p_j\bigr|\le K\delta_j\right)
   \ge\frac12 .
\end{equation}
\end{enumerate}
\end{lemma}
\begin{proof}
Throughout, $W$ is the contact envelope of \noderef{ContactEnv}, states are
written $z_n=(x_n,t_n)$, and $\delta_j=2^{-j}$.

\emph{Step 0: the constants.}
Let $C_f>0$ be the constant of \noderef{FitParamBound} for this $C_0$, let
$C_2>0$ be the constant of \noderef{DPLogVelocity} for this $C_0$ and
$\kappa$, and put
\[
  M=8,\qquad C_\Gamma=\max\{C_f,1\}>0 .
\]
Apply \noderef{StoppedWalkTravel} with the present standing data and this
$C_*$: it supplies an integer $T\ge8$ and a real $R\ge2$ such that every
stopped contact walk for $(\kappa,H,T,R,v,N,\epsilon)$ with $N\ge1$ satisfying
\eqref{eq:ml-contact-error-budget-hyp} obeys both of its conclusions, namely the
travel inequality in the portion generality and
\begin{equation}\label{eq:ml-hit}
  \mathbb P\bigl(t_\tau=\delta_N\bigr)\ge\frac34 .
\end{equation}
Fix these $T$, $R$ and put $B=R+T+2\ge12$, so $L(B)=1+\log B>1$. The constant
$K\ge2$ is fixed in Step 7; everything before that is proved for every
$K\ge2$ simultaneously, with constants not depending on $K$, $N$ or $j$.

Now let $N\ge1$, let $\epsilon$ and the walk be as in the statement, let
$(\sigma_j,p_j,a_j)$ be marks with fits, and fix $j$ with $1\le j\le N$. Put
$r=\delta_j$ and
\[
  \mathcal R_j=\{t_{\sigma_j}=r\},\qquad
  \epsilon'=\frac{\epsilon}{r},\qquad
  y_n=\frac{x_n-x_{\sigma_j}-(t_n-r)p_j}{r},\qquad
  s_n=\frac{t_n}{r}.
\]
By \noderef{IsMarkData}, clause 5, each $\mathbf1_{\{\sigma_j\le n\}}p_j$ is
$\mathcal G_n$-measurable, so $p_j=\sum_{n\ge0}\mathbf1_{\{\sigma_j=n\}}
 \mathbf1_{\{\sigma_j\le n\}}p_j$ is $\mathcal M$-measurable, and likewise
$a_j$; also $z_{\sigma_j}$ and $z_\tau$ are measurable, because $\sigma_j$ and
$\tau$ are stopping times and each $z_n$ is $\mathcal G_n$-measurable by
\noderef{IsContactWalk}, clauses 4 and 9. Hence $\mathcal R_j$ and the event in
\eqref{eq:contact-finite-localization} are measurable, and all expectations
below are of measurable functions.

\emph{Step 1: four grid facts, from \noderef{IsContactWalk}.}
By \noderef{IsContactWalk}, clause 2, $\epsilon=2^{-k}$ for some nonnegative
integer $k$ and $0<\epsilon\le\delta_N/4=2^{-N-2}$; hence $k\ge N+2$.

\emph{(a) Every $t_n$ lies on the grid $\epsilon\mathbb Z$, and
$\delta_i\in\epsilon\mathbb Z$ for $i\le N$.} Indeed
$\delta_i=2^{-i}=2^{k-i}\epsilon$ with $k-i\ge2$ a nonnegative integer, and
$1=2^k\epsilon$. For the states: $t_0=1$ by clause 5, and by clause 7 together
with \noderef{IsStepPair} --- whose two steps at a live state have time
coordinates $t_n+\epsilon$ and $t_n-\epsilon$ --- we have
\begin{equation}\label{eq:ml-timestep}
  t_{n+1}=t_n+\epsilon\xi_n\ \text{ on }\{n<\tau\},\qquad
  t_{n+1}=t_n\ \text{ on }\{n\ge\tau\},
\end{equation}
with $\xi_n\in\{1,-1\}$ by clause 3. Induction on $n$ gives
$t_n\in\epsilon\mathbb Z$. Consequently $S_n:=t_n/\epsilon$ is an integer for
every $n$, with $S_{n+1}=S_n+\xi_n$ on $\{n<\tau\}$.

\emph{(b) $t_{\sigma_j}\ge\delta_j$, with equality on $\{t_\tau=\delta_N\}$.}
Suppose first $t_{\sigma_j}<\delta_j$. Since $j\ge1$ we have
$t_0=1>\tfrac12\ge\delta_j$, so there is a least index $n_1\le\sigma_j$ with
$t_{n_1}\le\delta_j$, and $n_1\ge1$ with $t_{n_1-1}>\delta_j$. If
$n_1-1\ge\tau$ then $t_{n_1}=t_{n_1-1}>\delta_j$ by \eqref{eq:ml-timestep},
contradicting $t_{n_1}\le\delta_j$. So $n_1-1<\tau$ and
$t_{n_1}\ge t_{n_1-1}-\epsilon>\delta_j-\epsilon$ by
\eqref{eq:ml-timestep}; as $t_{n_1}$ and $\delta_j$ both lie on
$\epsilon\mathbb Z$ by (a), the two inequalities
$\delta_j-\epsilon<t_{n_1}\le\delta_j$ force $t_{n_1}=\delta_j$. Hence
$n_1<\sigma_j$ would contradict \noderef{IsMarkData}, clause 3, which gives
$t_n\ne\delta_j$ for $n<\sigma_j$; so $n_1=\sigma_j$ and
$t_{\sigma_j}=\delta_j$, contradicting $t_{\sigma_j}<\delta_j$. Therefore
$t_{\sigma_j}\ge\delta_j$.

On $\{t_\tau=\delta_N\}$ we have $t_\tau=\delta_N\le\delta_j$, since
$j\le N$. By \noderef{IsMarkData}, clause 4, either
$t_{\sigma_j}=\delta_j$ --- which is the assertion --- or
$\sigma_j=\tau$, in which case $t_{\sigma_j}=t_\tau\le\delta_j$ and the
inequality just proved gives $t_{\sigma_j}=\delta_j$ as well. So
\begin{equation}\label{eq:ml-Rj}
  \{t_\tau=\delta_N\}\subseteq\mathcal R_j .
\end{equation}

\emph{(c) $\epsilon'\le1/4$, and $\epsilon'$ is an inverse power of two
dividing $1$ and $M$.} Since $j\le N$ we have $\delta_N\le\delta_j=r$, so
$\epsilon'=\epsilon/r\le\epsilon/\delta_N\le1/4$; and
$\epsilon'=2^{-k}/2^{-j}=2^{j-k}$ with $k-j\ge2$, so $\epsilon'=2^{j-k}$
divides $1$ and divides the integer $M$.

\emph{(d) Before any step of the portion the primed time exceeds
$4\epsilon'$.} If $n<\tau$ then $t_n>\delta_N$ by \noderef{IsContactWalk},
clause 10, so
\begin{equation}\label{eq:ml-primed-time}
  s_n=\frac{t_n}{r}>\frac{\delta_N}{r}\ge\frac{4\epsilon}{r}=4\epsilon',
\end{equation}
using $\delta_N\ge4\epsilon$ from clause 2.

\emph{Step 2: the fit-parameter bound, i.e.\ (i).}
On $\mathcal R_j$, \noderef{IsMarkData}, clause 6, gives that $(p_j,a_j)$ is a
fit for $v$ at $z_{\sigma_j}=(x_{\sigma_j},r)$ of radius $r$ with constant
$C_0$, in the sense of \noderef{IsFit}. By \noderef{IsContactWalk}, clause 12,
$|x_{\sigma_j}|\le R+1\le B$, and $|r|=\delta_j\le1\le B$, while $B\ge2$ and
$r>0$. Hence \noderef{FitParamBound}, applied to $v$ with the centering point
$(0,1)$ --- the second alternative of its hypothesis --- the box size $K'=B$
and the radius $r$, gives
\[
  |p_j|+|a_j|\le C_f\bigl(L(B)+|\log r|\bigr)
   =C_f\bigl(L(B)+j\log2\bigr)\le C_f\bigl(L(B)+j\bigr)\le\Gamma_j
\]
on $\mathcal R_j$, using $\log2<1$ and $C_f\le C_\Gamma$. Off $\mathcal R_j$,
clause 7 of \noderef{IsMarkData} gives $p_j=0$ and $a_j=0$, so the bound holds
there too. Finally $L(B)>1$ and $j\ge1$ give
$\Gamma_j=C_\Gamma(L(B)+j)\ge1\cdot2>1$, so $\Gamma_j\ge1$. This is (i).

\emph{Step 3: the fitted frame.}
Fix $\omega\in\mathcal R_j$ and write $p_0=p_j(\omega)$, $a_0=a_j(\omega)$,
$x_0=x_{\sigma_j}(\omega)$, so that $(p_0,a_0)$ is a fit for $v$ at $(x_0,r)$
of radius $r$ and $r>0$. Let
\[
  \bar v_\omega(y,s)=\frac{v(x,t)-v(x_0,r)}{r}
     -p_0\cdot y-\Bigl(a_0+\frac{|p_0|^2}{2}\Bigr)(s-1),
  \qquad x=x_0+r\bigl(y+p_0(s-1)\bigr),\quad t=rs,
\]
be the contact frame of \noderef{FrameChange} based at $(x_0,r)$ with scale
$r$ and parameters $(p_0,a_0)$, preceded by the translation by $z_*=(0,1)$, and
let $\bar W_\omega$ be the contact envelope of \noderef{ContactEnv} formed from
$\bar v_\omega$ and the same $\kappa$. By \noderef{FrameEnvelopeCovariance},
whose hypotheses $1\le C_0$, $0\le\kappa$, $v\in\Aclass_{m,C_0}$, the Jensen
bound, $r>0$ and the fit are all available, its first three conclusions give
\begin{equation}\label{eq:ml-frame-class}
  \bar v_\omega\in\Aclass_{m,C_0},\quad
  \bar v_\omega\text{ is centered at }(0,1)\text{ with constant }C_0,\quad
  \bar v_\omega\text{ has the Jensen bound with the same }\kappa .
\end{equation}
These are exactly the standing hypotheses carried by $v$, so every statement
below that was quoted for a centered member of the class with a Jensen bound
applies to $\bar v_\omega$ and $\bar W_\omega$ with the same constants; this is
the only place the frame enters the constants, and it enters not at all.

Let $n$ satisfy $\sigma_j(\omega)\le n<\tau(\omega)$. Then $t_n>\delta_N\ge
4\epsilon>\epsilon$ by Step 1(d), the pair $(z^+_n,z^-_n)$ is an optimizing
step pair of size $\epsilon$ at $z_n$ for $v$ and $\kappa$ by
\noderef{IsContactWalk}, clause 6, and $\Penv_\epsilon W(\cdot,t_n)$ has
gradient $Q_0(n)$ at $x_n$ by clause 8. The images of $z_n$, $z^\pm_n$ under
the inverse frame map
$(x',t')\mapsto\bigl((x'-x_0-(t'-r)p_0)/r,\ t'/r\bigr)$ are
$\zeta_n=(y_n,s_n)$ and $\zeta^\pm_n$. So
\noderef{FrameVelocityCovariance} applies at this state and gives: the pair
$(\zeta^+_n,\zeta^-_n)$ is an optimizing step pair of size $\epsilon'$ at
$\zeta_n$ for $\bar v_\omega$ and $\kappa$; the primed velocities are
\begin{equation}\label{eq:ml-primed-velocities}
  Q'_\pm(n)=Q_\pm(n)-p_0,\qquad
  Q'_0(n)=Q_0(n)-p_0
   =\nabla_y\Penv_{\epsilon'}\bar W_\omega(y_n,s_n);
\end{equation}
and the discrepancies are unchanged,
\begin{equation}\label{eq:ml-same-discrepancies}
  \bigl|Q'_+(n)-Q'_0(n)\bigr|=\bigl|Q_+(n)-Q_0(n)\bigr|,\qquad
  \bigl|Q'_-(n)-Q'_0(n)\bigr|=\bigl|Q_-(n)-Q_0(n)\bigr| .
\end{equation}
No contact of $W$ with $v$ at the mark is used, and the fit and the frame are
determined by $\mathcal G_{\sigma_j}$-measurable data, hence fixed once the
mark is reached.

\emph{Step 4: the portion and its stopping index.}
Let $K\ge2$. Define $\tau'_j(\omega)$, for $\omega\in\mathcal R_j$, to be the
least $n\ge\sigma_j(\omega)$ with
\[
  n=\tau(\omega)\quad\text{or}\quad s_n(\omega)\ge M
  \quad\text{or}\quad |y_n(\omega)|\ge K,
\]
which exists because $n=\tau(\omega)\ge\sigma_j(\omega)$ qualifies; and put
$\tau'_j=\sigma_j$ off $\mathcal R_j$. Then
\begin{equation}\label{eq:ml-portion-order}
  \sigma_j\le\tau'_j\le\tau ,
\end{equation}
and for $\sigma_j\le n<\tau'_j$ we have $n<\tau$, $s_n<M$ and $|y_n|<K$.
Moreover $\tau'_j$ is a stopping time: $\{\tau'_j\le n\}$ is a finite union, over
$i\le n$, of intersections of $\{\sigma_j=i\}$ with events determined by
$z_0,\dots,z_n$ and by $\mathbf1_{\{\sigma_j\le i\}}p_j$, all of which lie in
$\mathcal G_n$ by \noderef{IsContactWalk}, clause 4, and
\noderef{IsMarkData}, clauses 1 and 5. Off $\mathcal R_j$ the portion is empty:
$\tau'_j=\sigma_j$, and there $p_j=0$, so $y_{\sigma_j}=0$ and every sum over
$\sigma_j\le n<\tau'_j$ is empty. On $\mathcal R_j$ we have $t_{\sigma_j}=r$,
hence $y_{\sigma_j}=0$ as well. In all cases
\begin{equation}\label{eq:ml-y-start}
  y_{\sigma_j}=0,\qquad
  y_n=\frac{x_n-x_{\sigma_j}-(t_n-t_{\sigma_j})p_j}{r}
  \quad(\sigma_j\le n\le\tau'_j),
\end{equation}
the second identity because on $\mathcal R_j$ we may replace $r$ by
$t_{\sigma_j}$ and off $\mathcal R_j$ we have $p_j=0$.

\emph{Step 5: the frame error budget.}
By \eqref{eq:ml-same-discrepancies}, for $\sigma_j\le n<\tau'_j$ the primed
squared discrepancy equals $\mathcal E_n$. Since
$\{n:\sigma_j\le n<\tau'_j\}\subseteq\{n:n<\tau\}$ by
\eqref{eq:ml-portion-order} and the summands are nonnegative, we get pointwise
$\sum_{\sigma_j\le n<\tau'_j}\mathcal E_n\le\sum_{n<\tau}\mathcal E_n$, and on
an empty portion the left side is $0$. Taking expectations and using
\eqref{eq:ml-contact-error-budget-hyp},
\begin{equation}\label{eq:contact-frame-error-budget}
  \mathbb E\sum_{\sigma_j\le n<\tau'_j}
    \bigl(|Q'_+(n)-Q'_0(n)|^2+|Q'_-(n)-Q'_0(n)|^2\bigr)
   =\mathbb E\sum_{\sigma_j\le n<\tau'_j}\mathcal E_n
   \le B_0:=C_*L(B)^4 .
\end{equation}
This is an unconditional estimate; it asserts nothing after conditioning on the
history at the mark.

\emph{Claim T$'$ (the conditional time walk).} Put $a=\delta_N/\epsilon$ and
$b=M/\epsilon'$, which are integers by Step 1(a) and (c), and note
$0\le a$. Let $h:\mathbb Z\to\R$ vanish outside $\{a,\dots,b\}$ and be
nonnegative there, and let $u:\mathbb Z\to[0,\infty)$ satisfy
\begin{equation}\label{eq:ml-poisson}
  \tfrac12\bigl(h(i+1)+h(i-1)\bigr)-h(i)=-u(i)\qquad(a<i<b).
\end{equation}
Then
\[
  \mathbb E\Bigl[\sum_{\sigma_j\le n<\tau'_j}u(S_n)\ \Big|\
      \mathcal G_{\sigma_j}\Bigr]\le h(S_{\sigma_j})\qquad\text{a.e.}
\]
Moreover, if $\ell:\mathbb Z\to\R$ is bounded and satisfies
$\tfrac12(\ell(i+1)+\ell(i-1))=\ell(i)$ for $a<i<b$, then
$\mathbb E[\ell(S_{\tau'_j})\mid\mathcal G_{\sigma_j}]=\ell(S_{\sigma_j})$
almost everywhere.

\emph{Proof of Claim T$'$.} Write $\rho=\sigma_j$ and $\varsigma=\tau'_j$. We
first record that
\begin{equation}\label{eq:ml-inside}
  a<S_n<b\qquad(\rho\le n<\varsigma).
\end{equation}
Indeed such an $n$ has $n<\tau$ by Step 4, so $t_n>\delta_N$ by
\noderef{IsContactWalk}, clause 10, i.e.\ $S_n>a$; and $s_n<M$ by Step 4, i.e.\
$S_n=s_n/\epsilon'<M/\epsilon'=b$. (If the portion is empty,
\eqref{eq:ml-inside} is vacuous.)

For $n\ge0$ put
\[
  Y_n=h\bigl(S_{(\rho\vee n)\wedge\varsigma}\bigr)
      +\sum_{\rho\le m<n\wedge\varsigma}u(S_m),
\]
a bounded function for each $n$, since $h$ and $u$ are bounded on the finite set
$\{a,\dots,b\}$ and vanish or are irrelevant elsewhere, and the sum has at most
$n$ terms.

Define
\[
  \mathcal G_\rho=\bigl\{A\subseteq\Omega:
     A\cap\{\rho\le n\}\in\mathcal G_n\ \text{for every }n\bigr\},
\]
the $\sigma$-algebra of events known at the stopping time $\rho$. It is a
$\sigma$-algebra: it contains $\Omega$ because $\{\rho\le n\}\in\mathcal G_n$,
$\rho$ being a stopping time; it is closed under complements because
$A^{\rm c}\cap\{\rho\le n\}=\{\rho\le n\}\setminus(A\cap\{\rho\le n\})$; and it
is closed under countable unions because intersection distributes over them. It
is contained in $\mathcal M$, since $A=\bigcup_n(A\cap\{\rho\le n\})$ and each
$A\cap\{\rho\le n\}\in\mathcal G_n\subseteq\mathcal M$. The only property of
$\mathcal G_\rho$ used below is the defining one.

Let $A\in\mathcal G_\rho$. We
claim $\int_AY_{n+1}\,d\mu=\int_AY_n\,d\mu$ for every $n$. Split $A$ into three
measurable pieces.

On $\{\rho>n\}$: then $\rho\ge n+1$, so
$(\rho\vee n)\wedge\varsigma=\rho\wedge\varsigma=\rho$ and
$(\rho\vee(n+1))\wedge\varsigma=\rho$ as well, using $\rho\le\varsigma$; both
sums are empty. Hence $Y_{n+1}=Y_n=h(S_\rho)$ there.

On $\{\rho\le n\}\cap\{\varsigma\le n\}$: then
$(\rho\vee n)\wedge\varsigma=\varsigma$ and
$(\rho\vee(n+1))\wedge\varsigma=\varsigma$, and both sums equal
$\sum_{\rho\le m<\varsigma}u(S_m)$. Hence $Y_{n+1}=Y_n$ there.

On $E_n:=A\cap\{\rho\le n<\varsigma\}$: here
$(\rho\vee n)\wedge\varsigma=n$ and
$(\rho\vee(n+1))\wedge\varsigma=n+1$, and the sum gains exactly the term
$u(S_n)$, so
\[
  Y_{n+1}-Y_n=h(S_{n+1})-h(S_n)+u(S_n).
\]
The event $E_n$ lies in $\mathcal G_n$: $A\cap\{\rho\le n\}\in\mathcal G_n$ by
the definition of $\mathcal G_\rho$, and
$\{n<\varsigma\}$ is the complement of $\{\varsigma\le n\}\in\mathcal G_n$,
$\varsigma$ being a stopping time by Step 4. On $E_n$ we have $n<\tau$, so
$S_{n+1}=S_n+\xi_n$ with $\xi_n\in\{1,-1\}$, and $a<S_n<b$ by
\eqref{eq:ml-inside}; writing
\[
  h(S_{n+1})=\frac{1+\xi_n}{2}h(S_n+1)+\frac{1-\xi_n}{2}h(S_n-1)
   =\frac{h(S_n+1)+h(S_n-1)}{2}
     +\frac{\xi_n}{2}\bigl(h(S_n+1)-h(S_n-1)\bigr),
\]
the coefficient of $\xi_n$ is bounded and $\mathcal G_n$-measurable, so the
pull-out property of the conditional expectation and
$\mathbb E[\xi_n\mid\mathcal G_n]=0$ (\noderef{IsContactWalk}, clause 3) give
$\int_{E_n}h(S_{n+1})\,d\mu
 =\int_{E_n}\tfrac12(h(S_n+1)+h(S_n-1))\,d\mu
 =\int_{E_n}(h(S_n)-u(S_n))\,d\mu$ by \eqref{eq:ml-poisson}. Hence
$\int_{E_n}(Y_{n+1}-Y_n)\,d\mu=0$, which proves the claim.

Therefore $\int_AY_n\,d\mu=\int_AY_0\,d\mu=\int_Ah(S_\rho)\,d\mu$ for every
$n$, since $Y_0=h(S_{\rho\wedge\varsigma})=h(S_\rho)$ and the sum in $Y_0$ is
empty. As $h\ge0$ on the range \eqref{eq:ml-inside} and $h$ vanishes elsewhere,
dropping the $h$ term gives
$\int_A\sum_{\rho\le m<n\wedge\varsigma}u(S_m)\,d\mu
 \le\int_Ah(S_\rho)\,d\mu$. The sums are nonnegative and nondecreasing in $n$,
with pointwise limit $X:=\sum_{\rho\le m<\varsigma}u(S_m)$ because $\varsigma$
is finite, so monotone convergence gives
\begin{equation}\label{eq:ml-condbound}
  \int_AX\,d\mu\le\int_Ah(S_\rho)\,d\mu\qquad(A\in\mathcal G_\rho).
\end{equation}
Taking $A=\Omega$ shows $X$ is integrable, since it is nonnegative, measurable
and has finite integral. Now both $\mathbb E[X\mid\mathcal G_\rho]$ and
$h(S_\rho)$ are $\mathcal G_\rho$-measurable and integrable, and
\eqref{eq:ml-condbound} says
$\int_A\mathbb E[X\mid\mathcal G_\rho]\,d\mu\le\int_Ah(S_\rho)\,d\mu$ for every
$A\in\mathcal G_\rho$; choosing
$A=\{\mathbb E[X\mid\mathcal G_\rho]>h(S_\rho)\}$ forces that event to be null,
so $\mathbb E[X\mid\mathcal G_\rho]\le h(S_\rho)$ a.e. For the last assertion,
run the same computation with $u=0$ and $h=\ell$: the three-piece argument used
neither $\ell\ge0$ nor the vanishing of $\ell$ outside $\{a,\dots,b\}$, only the
mean-value identity on $(a,b)$ and \eqref{eq:ml-inside}, so
$\int_A\ell(S_{(\rho\vee n)\wedge\varsigma})\,d\mu
 =\int_A\ell(S_\rho)\,d\mu$ for every $n$; since $\varsigma$ is finite,
$(\rho\vee n)\wedge\varsigma\to\varsigma$ pointwise, and $\ell$ is bounded while
$\mu$ is a probability measure by \noderef{IsContactWalk}, clause 1, so
dominated convergence gives
$\int_A\ell(S_\varsigma)\,d\mu=\int_A\ell(S_\rho)\,d\mu$ for every
$A\in\mathcal G_\rho$, i.e.\ the stated conditional identity. This proves
Claim T$'$.

\emph{Step 6: the conditional time and velocity budgets.}
We use three instances of Claim T$'$. Note first that
$\mathcal R_j\in\mathcal G_{\sigma_j}$ in the sense of Claim T$'$: for every
$n$,
\[
  \mathcal R_j\cap\{\sigma_j\le n\}
   =\bigcup_{i\le n}\bigl(\{\sigma_j=i\}\cap\{t_i=\delta_j\}\bigr),
\]
and each $\{\sigma_j=i\}$ lies in $\mathcal G_i\subseteq\mathcal G_n$ because
$\sigma_j$ is a stopping time by \noderef{IsMarkData}, clause 1, while
$\{t_i=\delta_j\}\in\mathcal G_i\subseteq\mathcal G_n$ because $z_i$ is
$\mathcal G_i$-measurable by \noderef{IsContactWalk}, clause 4. We also note
that $S_{\sigma_j}=r/\epsilon=1/\epsilon'$
on $\mathcal R_j$, while off $\mathcal R_j$ the portion is empty so every sum
below vanishes there; since $\mathcal R_j$ is $\mathcal G_{\sigma_j}$-measurable
we may therefore insert the indicator $\mathbf1_{\mathcal R_j}$ on the right of
each conditional bound.

\emph{(a) The number of steps.} Take $h_1(i)=(b-i)(i-a)$ on
$\{a,\dots,b\}$ and $0$ outside, and $u=\mathbf1_{(a,b)}$. Then $h_1\ge0$ on
$\{a,\dots,b\}$ and, for $a<i<b$, writing $h_1(i)=-i^2+(a+b)i-ab$, the second
difference of $-i^2$ is $-2$, so
$\tfrac12(h_1(i+1)+h_1(i-1))-h_1(i)=-1=-u(i)$. By \eqref{eq:ml-inside} the
summand $u(S_n)$ equals $1$ for every $n$ in the portion, so Claim T$'$ gives,
with $n'=\tau'_j-\sigma_j$ the number of steps of the portion and
$S_{\sigma_j}=1/\epsilon'$ on $\mathcal R_j$,
\[
  \mathbb E\bigl[n'\mid\mathcal G_{\sigma_j}\bigr]
   \le\mathbf1_{\mathcal R_j}\,h_1(1/\epsilon')
   =\mathbf1_{\mathcal R_j}\Bigl(\frac{M}{\epsilon'}-\frac1{\epsilon'}\Bigr)
      \Bigl(\frac1{\epsilon'}-a\Bigr)
   \le\mathbf1_{\mathcal R_j}\frac{M}{(\epsilon')^2},
\]
using $a\ge0$ and $M/\epsilon'-1/\epsilon'\le M/\epsilon'$. Hence
\begin{equation}\label{eq:ml-time-budget}
  \mathbb E\bigl[(\epsilon')^2n'\mid\mathcal G_{\sigma_j}\bigr]
   \le M\,\mathbf1_{\mathcal R_j}\qquad\text{a.e.}
\end{equation}

\emph{(b) Visit counts.} Fix an integer $i_0$ with $a<i_0<b$ and take
\[
  G_{i_0}(i)=\frac{2\min(i-a,\,i_0-a)\bigl(b-\max(i,i_0)\bigr)}{b-a}
\]
on $\{a,\dots,b\}$ and $0$ outside. Then $G_{i_0}\ge0$,
$G_{i_0}(a)=G_{i_0}(b)=0$, and $G_{i_0}$ is affine in $i$ on
$\{a,\dots,i_0\}$ and on $\{i_0,\dots,b\}$, so its second difference vanishes
for $a<i<b$, $i\ne i_0$, while at $i=i_0$, writing $q=b-a$,
$\alpha=i_0-a$,
\[
  \tfrac12\bigl(G_{i_0}(i_0-1)+G_{i_0}(i_0+1)\bigr)-G_{i_0}(i_0)
   =\frac{(\alpha-1)(q-\alpha)+\alpha(q-\alpha-1)}{q}
      -\frac{2\alpha(q-\alpha)}{q}=-1,
\]
because $(\alpha-1)(q-\alpha)+\alpha(q-\alpha-1)=2\alpha(q-\alpha)-q$. So
\eqref{eq:ml-poisson} holds with $u=\mathbf1_{\{i_0\}}$, and Claim T$'$ gives,
with $V_{i_0}=\#\{n:\sigma_j\le n<\tau'_j,\ S_n=i_0\}$,
\begin{equation}\label{eq:ml-visits}
  \mathbb E\bigl[V_{i_0}\mid\mathcal G_{\sigma_j}\bigr]
   \le\mathbf1_{\mathcal R_j}G_{i_0}(1/\epsilon')
   \le\mathbf1_{\mathcal R_j}\frac{b-a}{2}
   \le\mathbf1_{\mathcal R_j}\frac{M}{2\epsilon'},
\end{equation}
the middle step because $\min(i-a,\alpha)(b-\max(i,i_0))\le\alpha(q-\alpha)\le
q^2/4$ for every $i$, and the last because $b-a\le b=M/\epsilon'$.

\emph{(c) The velocity budget.} Let $\sigma_j\le n<\tau'_j$ and
$\omega\in\mathcal R_j$. Then $|y_n|<K\le K+M+2$ and $0<s_n<M\le K+M+2$ by
Step 4, $0<\epsilon'\le1/4\le1$ by Step 1(c), $\epsilon'<s_n$ by
\eqref{eq:ml-primed-time}, and $K+M+2\ge2$. By \eqref{eq:ml-frame-class} the
function $\bar v_\omega$ is a member of $\Aclass_{m,C_0}$ centered at $(0,1)$
with the Jensen bound at the same $\kappa$, and by
\eqref{eq:ml-primed-velocities} the vector $Q'_0(n)$ is the gradient of
$\Penv_{\epsilon'}\bar W_\omega(\cdot,s_n)$ at $y_n$. Hence
\noderef{DPLogVelocity}, applied to $\bar v_\omega$ with the box size
$K+M+2$, gives
\begin{equation}\label{eq:ml-Q0prime}
  |Q'_0(n)|\le C_2\bigl(L(K+M+2)+|\log s_n|\bigr),
\end{equation}
with $C_2$ the constant of \noderef{DPLogVelocity}, which depends only on
$C_0$, $\dim F$ and $\kappa$ --- in particular not on the frame, which is why
the bound is uniform in $j$, $N$ and $\omega$.

From $|\alpha+\beta|^2\le2|\alpha|^2+2|\beta|^2$ and
\eqref{eq:ml-Q0prime},
$(\epsilon')^2|Q'_0(n)|^2\le2C_2^2(\epsilon')^2L(K+M+2)^2
  +2C_2^2(\epsilon')^2|\log s_n|^2$. Summing over the portion, the first part
contributes $2C_2^2L(K+M+2)^2(\epsilon')^2n'$, whose conditional expectation is
at most $2C_2^2L(K+M+2)^2M\mathbf1_{\mathcal R_j}$ by
\eqref{eq:ml-time-budget}. For the second, $s_n=S_n\epsilon'$ and
$a<S_n<b$, so grouping the summands by the level $i_0=S_n$ and using
\eqref{eq:ml-visits},
\[
  \mathbb E\Bigl[\sum_{\sigma_j\le n<\tau'_j}
     (\epsilon')^2|\log s_n|^2\ \Big|\ \mathcal G_{\sigma_j}\Bigr]
   =(\epsilon')^2\sum_{i_0=a+1}^{b-1}\bigl|\log(i_0\epsilon')\bigr|^2
       \mathbb E\bigl[V_{i_0}\mid\mathcal G_{\sigma_j}\bigr]
   \le\mathbf1_{\mathcal R_j}\frac{M\epsilon'}{2}
       \sum_{i_0=a+1}^{b-1}\bigl|\log(i_0\epsilon')\bigr|^2 .
\]
The levels $s=i_0\epsilon'$ with $a<i_0<b$ lie in $(0,M)$. Split the sum. For
$i_0\epsilon'\le1$ the function $s\mapsto|\log s|^2$ is nonnegative and
decreasing on $(0,1]$, so its right-endpoint grid sum obeys
$\epsilon'\sum_{i_0\epsilon'\le1}|\log(i_0\epsilon')|^2
 \le\int_0^1|\log s|^2\,ds=2$, the value of the integral following from the
substitution $s=e^{-\upsilon}$, which turns it into
$\int_0^\infty\upsilon^2e^{-\upsilon}\,d\upsilon=2$. For
$1<i_0\epsilon'<M$ we have $|\log(i_0\epsilon')|^2\le(\log M)^2$ and there are
at most $b=M/\epsilon'$ such $i_0$, so
$\epsilon'\sum_{i_0\epsilon'>1}|\log(i_0\epsilon')|^2\le M(\log M)^2$.
Therefore, using $L(K+M+2)\ge1$,
\begin{equation}\label{eq:contact-frame-time-budget}
  \mathbb E\Bigl[\sum_{\sigma_j\le n<\tau'_j}
     (\epsilon')^2|Q'_0(n)|^2\ \Big|\ \mathcal G_{\sigma_j}\Bigr]
   \le C_ML(K+M+2)^2\,\mathbf1_{\mathcal R_j},\qquad
  C_M=2C_2^2\Bigl(M+\frac M2\bigl(2+M(\log M)^2\bigr)\Bigr).
\end{equation}
Integrating \eqref{eq:ml-time-budget} and
\eqref{eq:contact-frame-time-budget} over $\Omega$ removes the indicators and
gives the unconditional bounds
\begin{equation}\label{eq:ml-integrated}
  \mathbb E\bigl[(\epsilon')^2n'\bigr]\le M,\qquad
  \mathbb E\sum_{\sigma_j\le n<\tau'_j}(\epsilon')^2|Q'_0(n)|^2
    \le C_ML(K+M+2)^2 .
\end{equation}

The first of these two unconditional bounds is exactly
\noderef{ExpectedPortionLengthBound}, applied to the lattice data of Claim
T$'$ --- the process $S$, the indices $\sigma_j\le\tau'_j\le\tau$, the window
$(a,b)$ and the whole-space instance of the comparison --- with the test pair
$h_1$, $u$ of (a); the conditional form \eqref{eq:ml-time-budget} is kept here
because the velocity estimate above reads it before the integration.

\emph{Step 7: the frame travel bound and the choice of $K$.}
Apply the travel inequality of \noderef{StoppedWalkTravel}, available by
Step 0, with the stopping times $\rho=\sigma_j$ and $\varsigma=\tau'_j$, which
satisfy $\sigma_j\le\tau'_j\le\tau$ by \eqref{eq:ml-portion-order}, with
$p=p_j$, $C_p=\Gamma_j$ --- admissible by Step 2 --- and with $r=\delta_j>0$;
the required measurability of $\mathbf1_{\{\sigma_j\le n\}}p_j$ is
\noderef{IsMarkData}, clause 5. By \eqref{eq:ml-y-start} its left-hand side is
exactly $\mathbb E\sup_{\sigma_j\le n\le\tau'_j}|y_n|$, and by
\eqref{eq:ml-primed-velocities} its first right-hand term involves
$|Q_0(n)-p_j|=|Q'_0(n)|$ with the factor $(\epsilon/r)^2=(\epsilon')^2$. So the
clause asserts that $\sup_{\sigma_j\le n\le\tau'_j}|y_n|$ is integrable and
\[
  \mathbb E\sup_{\sigma_j\le n\le\tau'_j}|y_n|
   \le2\left(\mathbb E\sum_{\sigma_j\le n<\tau'_j}
        (\epsilon')^2|Q'_0(n)|^2\right)^{1/2}
     +\left(\mathbb E(\epsilon')^2(\tau'_j-\sigma_j)\right)^{1/2}
      \left(\mathbb E\sum_{\sigma_j\le n<\tau'_j}\mathcal E_n\right)^{1/2}.
\]
Inserting \eqref{eq:ml-integrated} and
\eqref{eq:contact-frame-error-budget},
\begin{equation}\label{eq:contact-frame-travel}
  \mathbb E\sup_{\sigma_j\le n\le\tau'_j}|y_n|
   \le2\sqrt{C_M}\,L(K+M+2)+\sqrt{M}\sqrt{B_0}
   \le C'_M\bigl(L(K+M+2)+\sqrt{B_0}\bigr),
\end{equation}
with $C'_M=\max\{2\sqrt{C_M},\sqrt M\}$, a constant depending only on $M$,
$C_0$, $\kappa$ and $\dim F$ --- and not on $K$, $N$ or $j$.

Since $B_0=C_*L(B)^4$ and $M$ are now fixed and
$L(K+M+2)/K\to0$ as $K\to\infty$, we may fix once and for all
\begin{equation}\label{eq:ml-choiceK}
  K\ge2\quad\text{with}\quad
  \frac{C'_M\bigl(L(K+M+2)+\sqrt{B_0}\bigr)}{K}\le\frac18 .
\end{equation}
This choice uses only $M$, $B_0$ and the fixed data, so it is uniform in $N$
and $j$, as required by the order of the quantifiers in the statement.

\emph{Step 8: the two cutoff probabilities.}
Put
\[
  \mathcal B^1_j=\mathcal R_j\cap\bigl\{|y_{\tau'_j}|\ge K\bigr\},\qquad
  \mathcal B^2_j=\mathcal R_j\cap\bigl\{s_{\tau'_j}\ge M\bigr\}.
\]
On $\mathcal B^1_j$ we have
$|y_{\tau'_j}|\le\sup_{\sigma_j\le n\le\tau'_j}|y_n|$, so the pointwise
inequality $K\mathbf1_{\mathcal B^1_j}\le\sup_{\sigma_j\le n\le\tau'_j}|y_n|$
holds; integrating it --- legitimate because that supremum is integrable by
Step 7 --- and using \eqref{eq:contact-frame-travel} and
\eqref{eq:ml-choiceK},
\begin{equation}\label{eq:ml-prob1}
  \mathbb P\bigl(\mathcal B^1_j\bigr)
   \le\frac{\mathbb E\sup_{\sigma_j\le n\le\tau'_j}|y_n|}{K}
   \le\frac{C'_M(L(K+M+2)+\sqrt{B_0})}{K}\le\frac18 .
\end{equation}
For the other cutoff take $\ell(i)=(i-a)/(b-a)$ for $a\le i\le b$,
$\ell(i)=0$ for $i<a$ and $\ell(i)=1$ for $i>b$. It is bounded and affine on
$\{a,\dots,b\}$, hence satisfies the mean-value identity for $a<i<b$, so the
last assertion of Claim T$'$ applies:
$\mathbb E[\ell(S_{\tau'_j})\mid\mathcal G_{\sigma_j}]=\ell(S_{\sigma_j})$
a.e. On $\mathcal R_j$ we have $S_{\sigma_j}=1/\epsilon'$ and $b=M/\epsilon'$,
so, with $a\ge0$ and $M\ge1$,
\[
  \ell(S_{\sigma_j})
   =\frac{1/\epsilon'-a}{M/\epsilon'-a}\le\frac1M ,
\]
the inequality because $M(1/\epsilon'-a)\le M/\epsilon'-a$ reduces to
$Ma\ge a$. Also $s_{\tau'_j}\ge M$ means $S_{\tau'_j}\ge b$, where
$\ell=1$, so $\mathbf1_{\mathcal B^2_j}\le\mathbf1_{\mathcal R_j}
 \ell(S_{\tau'_j})$ pointwise, $\ell$ being nonnegative. Since
$\mathcal R_j\in\mathcal G_{\sigma_j}$,
\begin{equation}\label{eq:ml-prob2}
  \mathbb P\bigl(\mathcal B^2_j\bigr)
   \le\mathbb E\bigl[\mathbf1_{\mathcal R_j}\ell(S_{\tau'_j})\bigr]
   =\mathbb E\Bigl[\mathbf1_{\mathcal R_j}
       \mathbb E\bigl[\ell(S_{\tau'_j})\mid\mathcal G_{\sigma_j}\bigr]\Bigr]
   =\mathbb E\bigl[\mathbf1_{\mathcal R_j}\ell(S_{\sigma_j})\bigr]
   \le\frac1M=\frac18 ,
\end{equation}
the middle equality by the characterizing property of the conditional
expectation applied to the set $\mathcal R_j\in\mathcal G_{\sigma_j}$.

\emph{Step 9: the combination, i.e.\ (ii).}
Let $\mathcal G=\{t_\tau=\delta_N\}$, so $\mathbb P(\mathcal G)\ge3/4$ by
\eqref{eq:ml-hit} and $\mathcal G\subseteq\mathcal R_j$ by
\eqref{eq:ml-Rj}. Let
$\omega\in\mathcal G\setminus(\mathcal B^1_j\cup\mathcal B^2_j)$. Then
$\omega\in\mathcal R_j$, and since $\omega\notin\mathcal B^1_j\cup\mathcal
B^2_j$ we have $|y_{\tau'_j}|<K$ and $s_{\tau'_j}<M$. By the definition of
$\tau'_j$ in Step 4, at least one of the three alternatives holds at
$n=\tau'_j$; the second and third have just been excluded, so
$\tau'_j(\omega)=\tau(\omega)$. Hence $|y_\tau|<K$, which by
\eqref{eq:ml-y-start} and $t_{\sigma_j}=\delta_j$ on $\mathcal R_j$ reads
\[
  \bigl|x_\tau-x_{\sigma_j}-(t_\tau-\delta_j)p_j\bigr|<K\delta_j .
\]
Together with $t_\tau=\delta_N$ this places $\omega$ in the event of
\eqref{eq:contact-finite-localization}. Therefore that event contains
$\mathcal G\setminus(\mathcal B^1_j\cup\mathcal B^2_j)$ and, by
\eqref{eq:ml-hit}, \eqref{eq:ml-prob1} and \eqref{eq:ml-prob2},
\[
  \mathbb P\Bigl(t_\tau=\delta_N\ \text{and}\
     \bigl|x_\tau-x_{\sigma_j}-(t_\tau-\delta_j)p_j\bigr|\le K\delta_j\Bigr)
   \ge\mathbb P(\mathcal G)-\mathbb P(\mathcal B^1_j)
      -\mathbb P(\mathcal B^2_j)
   \ge\frac34-\frac18-\frac18=\frac12 .
\]

The event $\mathcal L_j$ is measurable, which is the remaining half of (ii).
Both $\tau$ and $\sigma_j$ take values in $\mathbb N$, so
\[
  \mathcal L_j=\bigcup_{n\ge0}\ \bigcup_{m\ge0}
    \bigl(\{\tau=n\}\cap\{\sigma_j=m\}\cap A_{n,m}\bigr),\qquad
  A_{n,m}=\{t_n=\delta_N\}\cap
    \bigl\{\bigl|x_n-x_m-(t_n-\delta_j)p_j\bigr|\le K\delta_j\bigr\},
\]
a countable union. Here $\{\tau=n\}\in\mathcal G_n\subseteq\mathcal M$ because
$\tau$ is a stopping time for $(\mathcal G_n)$ by clause 9 of
\noderef{IsContactWalk}, and $\{\sigma_j=m\}\in\mathcal G_m\subseteq\mathcal M$
because $\sigma_j$ is a stopping time by clause 1 of \noderef{IsMarkData}. The
maps $z_n$ and $z_m$ are $\mathcal G_n$- and $\mathcal G_m$-measurable by
clause 4 of \noderef{IsContactWalk}, the coordinate projections of
$\ell^2(F\times\R)$ are continuous and hence Borel, and on
$\{\sigma_j=m\}\cap\{\sigma_j\le m\}=\{\sigma_j=m\}$ the function $p_j$ agrees
with $\mathbf1_{\{\sigma_j\le m\}}p_j$, which is $\mathcal G_m$-measurable by
clause 5 of \noderef{IsMarkData}; so
$\{\sigma_j=m\}\cap A_{n,m}\in\mathcal G_{\max(n,m)}\subseteq\mathcal M$, the
norm and the linear operations being continuous. A countable union of
measurable sets is measurable, so $\mathcal L_j\in\mathcal M$; the same
decomposition applied to $\mathcal B^1_j$, $\mathcal B^2_j$ and $\mathcal G$
is what makes the probabilities used above meaningful.

This is \eqref{eq:contact-finite-localization}, and with Step 2 it completes
the proof.
\end{proof}
